\section{Convergence of the Gaussian resolvent trace}
\label{sec:trace}

The form convergence in Section~\ref{sec:microscopic} identifies each
fixed edge coordinate.  The autocorrelation, however, sums over every
coordinate, and the static variances are not summable.  We now prove the
additional trace estimate that permits this sum to pass to the limit.
Its mechanism is angular integration by parts: an edge gap of order
$a^{2/3}$ gives a factor $a^{-2/3}$ in the dual norm of the $a$th source.
The squares of these factors are summable.
The variational passage from forms to individual resolvents is the
standard one in the theory of varying Hilbert spaces
\cite[Theorem~2.4]{KuwaeShioya2003}; the summation over sources requires
the additional estimates below.

Throughout this section the disorder is Gaussian.  Write
\begin{equation}
 B_g(z):=\sum_{a\geq1}
 \big\langle x_a,(z+c_*H_\Pi)^{-1}x_a\big\rangle_\Pi,
 \qquad z>0.
 \label{eq:source-40}
\end{equation}
The sum is initially allowed to be infinite.  Recall that
$B_n(z)=\sum_{a=1}^n\langle x_a,(z+TH)^{-1}x_a\rangle_\mu$ and
$T=n^{2/3}$.

\begin{proposition}[Gaussian trace convergence]
\label{prop:trace-goe}
For almost every edge configuration $g$, the series defining $B_g(z)$
is finite for every $z>0$.  For every finite collection $z_1,\ldots,z_k>0$,
\[
 (B_n(z_1),\ldots,B_n(z_k))
 \ \Longrightarrow\ (B_g(z_1),\ldots,B_g(z_k)).
\]
Moreover, for each $z,\epsilon>0$,
\[
 \lim_{\ell\to\infty}\limsup_{n\to\infty}
 \mathbb P\left\{
 \sum_{a>\ell}\langle x_a,(z+TH)^{-1}x_a\rangle_\mu
 >\epsilon\right\}=0.
\]
\end{proposition}

We separate fixed-coordinate convergence from the estimate uniform in
the number of coordinates.  For a vector-valued function $f$, let
\[
 \|f\|_{z,n}^2=z\mu|f|^2+T\mathcal E(f),
\]
where the energy includes the sum over vector components.
The notation $\|f\|_\infty$ means $\sup_X|f(X)|$ in the Euclidean norm
of the target space, or its Hilbert--Schmidt norm for matrix targets.
We allow a polynomial number of components: summing an exponentially
small scalar remainder over that many components still gives a
negligible error for each fixed polynomial supremum bound.
The resolvent identity in variational form is
\[
 \langle S,(z+TH)^{-1}S\rangle_\mu
 =\sup_f\{2\mu[S\cdot f]-\|f\|_{z,n}^2\}
 =\left(\sup_{\|f\|_{z,n}\leq1}\mu[S\cdot f]\right)^2.
\]
This formula applies to both scalar and vector sources.

\begin{lemma}[Fixed coordinates]
\label{lem:trace-fixed}
For every fixed $\ell$ and every finite collection of positive $z$,
the corresponding sums over $1\leq a\leq\ell$ converge jointly to
the same restrictions of \eqref{eq:source-40}.
\end{lemma}

\begin{proof}
Use the quenched subsequential realization from
Section~\ref{sec:edge-law}.  For a fixed coordinate the source has
bounded $L^2$ norm on the represented sequence, by the static moment
convergence.  Its resolvent optimizer therefore has bounded $L^2$ norm
and bounded scaled energy.  Markov contraction also gives a polynomial
supremum bound: $\|x_a\|_\infty\leq n^{1/6}$ and
\[
 \|(z+TH)^{-1}x_a\|_\infty\leq z^{-1}n^{1/6}.
\]
The lower form bound \eqref{eq:source-38}, together with weak lower
semicontinuity of the mass term, bounds the upper limit of the
variational expression.  Source pairings pass to the limit by first
cutting off $x_a$, using static cylinder convergence, and then removing
the cutoff with the fixed-coordinate moment bounds.

Conversely, the correctors \eqref{eq:source-39} give the variational
lower bound for every smooth compact cylinder test: first take
$n\to\infty$ at a fixed recovery polynomial $Q_j$, and then take
$j\to\infty$.  The polynomial supremum exponent may depend on $j$,
which is harmless in this order of limits.  The cylinder core
in Proposition~\ref{prop:gradient-domain} extends it to the full form
domain.  The two variational bounds identify the limiting resolvent.
The same argument uses one subsequence for the finitely many sources
and parameters under consideration, proving the joint assertion.
\end{proof}

\subsection{Moving a vector pairing to the smallest cap}

For a vector test $f=(f_a)_{a>\ell}$, put
$F_p=\mathbb E[f\mid h_p]$, componentwise.  We use the ordinary
observation law $Q_\mu$, and all its expectations below are at fixed
disorder.  The scale is $r=n^{-1/3+\delta}$, with the fixed $\delta$
chosen in Section~\ref{sec:edge-law}.

\begin{lemma}[Observation replacement]
\label{lem:trace-replacement}
Uniformly over vector tests satisfying $\|f\|_{z,n}\leq1$ and a fixed
polynomial supremum bound,
\begin{equation}
 \sum_{a>\ell}\mu[f_ax_a]
 =n^{-1/3}\sum_{a>\ell}
 Q_\mu\big[(F_r)_a(U^tm_r)_a\big]+o_{\mathbb P}(1).
 \label{eq:source-41}
\end{equation}
The error is uniform in $\ell$ and in the number of vector components.
\end{lemma}

\begin{proof}
The posterior means of $f$ and $X$ are martingales as the observation
precision increases.  If $c<p\leq2c$, orthogonality of martingale
increments bounds the change in their pairing by
\[
 n^{-1/3}
 \big(Q_\mu|F_p-F_c|^2\big)^{1/2}
 \big(Q_\mu|m_p-m_c|^2\big)^{1/2}
 \leq Cn^{-2/3}c^{-1}
 \big(Q_\mu|m_p|^2\big)^{1/2}+o(1).
\]
Here \eqref{eq:source-16} bounds the first increment by
$Cn^{-1/3}c^{-1}$; its exponentially small remainder can be summed
over at most $n$ components with polynomially bounded supremum.
The moment estimate \eqref{eq:source-11}, with $M_p=(np)^{1/2}$,
makes the last bound of order
$n^{-1/6}p^{-1/2}$, apart from tight constants.
On a geometric list beginning at $r$, these deterministic weights sum
to $O(n^{-\delta/2})$.

For completeness, no simultaneous bound on the maximum of all these
random constants is required.  Restrict first to the common disorder
events on which \eqref{eq:source-16} holds with a fixed constant.
Average the weighted sum over the biased rotation law, use
\eqref{eq:source-11}, and then use $L_\circ\geq1/2$ and Markov's
inequality to return to ordinary disorder probability.  The probability
of the discarded disorder events can be made arbitrarily small.

At the fixed upper endpoint $p_*$, the difference from the original
pairing is bounded directly using \eqref{eq:source-16} and
$|X|=\sqrt n$.  It is $O(n^{-1/6})$.  Combining this endpoint estimate
with the geometric sum proves the lemma.  All inequalities use vector
norms, so their constants do not grow with the vector dimension.
\end{proof}

\subsection{Angular integration by parts and a summable tail}

We next bound the right side of \eqref{eq:source-41}.  By
\eqref{eq:source-12}, replacing $m_r$ there by $m_r^{\rm sp}$ costs
$o_{\mathbb P}(1)$, uniformly over the tests in the lemma.  Indeed its
annealed squared error is bounded by a constant times
$n^{-2/3}M_r^2n^{-1/20}=n^{\delta-1/20}$.
On the cap annulus, radial symmetry of the spherical posterior gives
\[
 n^{-1/3}(U^tm_r^{\rm sp})_a=t(\|h_r\|)u_a,
 \qquad u=u_r,
\]
on the active coordinate space, and zero on its orthogonal complement.
The saddle bounds make $t$ bounded.  Insert a smooth radial cutoff
supported in the annulus and equal to one on a smaller annulus containing
the event in \eqref{eq:source-7}.  The discarded contribution is
negligible by its power-exponential probability and the polynomial
bounds.  Write $\theta(u)$ for the resulting bounded radial multiplier.

Fix $m_0=50$ and define
$R_0=\sum_{j=1}^{m_0}u_j^2$.  On the tight spectral sets and for
sufficiently large $n$, these leading coordinates are all active.
For $a>\ell$ active and $j\leq m_0$, let $\delta_{aj}$ be the coefficient
in \eqref{eq:source-31}.  Taking $\ell$ beyond a tight index cutoff gives
$\delta_{aj}\geq c a^{2/3}$.  Define
\begin{equation}
 \begin{split}
 \mathcal V_a
 &=\sum_{j=1}^{m_0}\frac{u_j}{R_0\delta_{aj}}
       (u_j\partial_a-u_a\partial_j),\\
 |\mathcal V_a|
 &\leq Ca^{-2/3}\left(1+\frac{|u_a|}{\sqrt{R_0}}\right),\qquad
 |\operatorname{div}\mathcal V_a|
 \leq Ca^{-2/3}\frac{|u_a|}{R_0}.
 \end{split}
 \label{eq:source-42}
\end{equation}
The vector field is tangent to Euclidean spheres.  It therefore
annihilates $\theta$ and every radial term in the spherical log density.
The normalization of $\delta_{aj}$ gives the exact identity
\[
 \mathcal V_a\left(-\frac12h_r^tB_r^{-1}h_r\right)=-u_a.
\]

We record the integrability needed to use these fields.  In the biased
law $\mathbf P$, the first $m_0$ signal coordinates have a uniformly
bounded joint density on each tight spectral set, by the conditioning
construction of $\Pi$ and its finite-dimensional version.  Their signal
multipliers $B_{r,j}/r^2$ are bounded away from zero.  Adding independent
observation noise preserves a uniform density bound.  Hence, for every
$q<m_0/2$,
\[
 \mathbf E R_0^{-q}\leq C_q.
\]
For example, integrate the bounded density over the unit ball in
$\mathbb R^{m_0}$ and use $R_0^{-q}\leq1$ outside that ball.
All individual $u_a$ also have uniformly bounded moments of every
fixed order.  Thus, with
$s_\ell=\sum_{a>\ell}a^{-4/3}$ and with sums restricted to active indices,
\[
 \mathbf E\sum_{a>\ell}|\operatorname{div}\mathcal V_a|^2
 \leq Cs_\ell,
 \qquad
 \left\|\sum_{a>\ell}|\mathcal V_a|^2\right\|_{L^2(\mathbf P)}
 \leq Cs_\ell.
\]
The second estimate follows from Minkowski's inequality and
\[
 \left(\mathbf E\frac{u_a^4}{R_0^2}\right)^{1/2}
 \leq(\mathbf E u_a^8)^{1/4}
       (\mathbf E R_0^{-4})^{1/4}.
\]
In particular, the dimensions of the active cap never enter these
bounds.

The ordinary observation density in $u$ is proportional to
\[
 L_r(u)\exp\left\{\phi_r^{\rm sp}(h_r)
                         -\frac12h_r^tB_r^{-1}h_r\right\}.
\]
Integration by parts therefore gives
\[
 \begin{split}
 \sum_{a>\ell}Q_\mu[\theta(F_r)_a u_a]
 ={}&\sum_{a>\ell}Q_\mu[\theta\mathcal V_a(F_r)_a]\\
 &+\sum_{a>\ell}Q_\mu[\theta(F_r)_a\operatorname{div}\mathcal V_a]\\
 &+\sum_{a>\ell}Q_\mu[\theta(F_r)_a\mathcal V_a\log L_r].
 \end{split}
\]
These three terms correspond respectively to differentiating the test,
the vector field, and the comparison density.  The first is bounded by
Cauchy--Schwarz using
$Q_\mu\sum_a|\mathcal V_a|^2$ and
$Q_\mu\sum_a|\nabla_u(F_r)_a|^2$; the latter is bounded by the vector
version of \eqref{eq:source-26}.  The second uses
$Q_\mu|F_r|^2\leq z^{-1}$ and the divergence estimate above.
For the third, put $W_\ell=\sum_{a>\ell}|\mathcal V_a|^2$.  Then
\[
 \mathbf E\sum_{a>\ell}|\mathcal V_a\log L_r|^2
 \leq \|W_\ell\|_{L^2(\mathbf P)}
       \big(\mathbf E|\nabla_u\log L_r|^4\big)^{1/2}
 =o(1)s_\ell
\]
by \eqref{eq:source-12}.  Transferring these nonnegative bounds to
ordinary disorder with $L_\circ\geq1/2$ and Markov's inequality shows
that the entire pairing is
$O_{\mathbb P}(s_\ell^{1/2})+o_{\mathbb P}(1)$, uniformly over the tests.
More precisely, for each $\eta>0$ one first chooses a deterministic
index $L_\eta$ and bounds on the spectral, variance, and gradient
constants, outside disorder events of probability at most $\eta$.
The displayed moment estimates then hold with a constant $C_\eta$
for every $\ell\geq L_\eta$, uniformly in $n$.  This is the uniformity
in the tail index that will be used below.
The integration by parts is justified first with a cutoff away from
$R_0=0$ and then without it.  If $\chi(R_0/\varepsilon)$ is such a
cutoff, then
\[
 |\mathcal V_a\chi(R_0/\varepsilon)|
 \leq Ca^{-2/3}\frac{|u_a|}{R_0}
             \mathbf1_{\{R_0\leq2\varepsilon\}}.
\]
The squared norms of these cutoff terms sum to a quantity tending to
zero by the coordinate and inverse moments above.  They are therefore
harmless against the vector $F_r$ by Cauchy--Schwarz.  At fixed
$n$ the observation density is smooth on its active space.

\begin{proof}[Proof of Proposition~\ref{prop:trace-goe}]
Let $B_{n,>\ell}(z)$ denote the tail trace in the proposition.  If it
is at least $\epsilon$, its normalized optimizer is
\[
 f=\frac{(z+TH)^{-1}(x_a)_{a>\ell}}{\sqrt{B_{n,>\ell}(z)}}.
\]
It has $\|f\|_{z,n}=1$ and
$\|f\|_\infty\leq n^{1/6}/(z\sqrt\epsilon)$, since
$|(x_a)_{a>\ell}|\leq n^{1/6}$ pointwise.  It is therefore an admissible
test for the uniform estimates above.  Lemma~\ref{lem:trace-replacement}
and the angular calculation imply that its pairing with the source
vanishes in probability in the double limit $n\to\infty$, then
$\ell\to\infty$.  That pairing equals $\sqrt{B_{n,>\ell}(z)}$.
In the localized form just described, Markov's inequality gives, for
each fixed $\epsilon>0$ and every $\ell\geq L_\eta$,
\[
 \limsup_{n\to\infty}
 \mathbb P\{B_{n,>\ell}(z)>\epsilon\}
 \leq\eta+C_\eta\epsilon^{-1}s_\ell.
\]
Letting $\ell\to\infty$ and then $\eta\downarrow0$ proves the asserted
tail estimate.

Combine it with Lemma~\ref{lem:trace-fixed}.  The tail estimate and
tightness of one sufficiently long fixed-coordinate sum first give
tightness of the full $B_n(z)$.  Each limiting finite partial sum is
therefore bounded by the same tightness bounds, uniformly in its
length: apply convergence of that partial sum and its pointwise bound
by $B_n(z)$.  Monotonicity now shows that the limiting series is finite
almost surely.  The tail estimate then proves full
trace convergence, jointly at finitely many parameters.  Finiteness at
all positive $z$ follows first at positive rationals and then everywhere
by resolvent monotonicity.
\end{proof}

\subsection{A quadratic feature needed for disorder comparison}

One more Gaussian estimate will allow the trace to be compared across
entry distributions.  Regard $XX^t/n$ as a vector in the space of
matrices with the Hilbert--Schmidt inner product, and set
\[
 B_n^{[2]}(z)=\sum_{i,j=1}^n
 \left\langle\frac{X_iX_j}{n},
 (z+TH)^{-1}\frac{X_iX_j}{n}\right\rangle_\mu.
\]
This feature has pointwise squared norm one, so
$0\leq B_n^{[2]}(z)\leq z^{-1}$.

\begin{proposition}[Vanishing quadratic-feature trace]
\label{prop:trace-second}
For Gaussian disorder and every $z>0$,
\begin{equation}
 B_n^{[2]}(z)\longrightarrow0
 \quad\text{in probability}.
 \label{eq:source-43}
\end{equation}
\end{proposition}

\begin{proof}
Test the feature against a vector $f$ with $\|f\|_{z,n}\leq1$ and a
polynomial supremum bound.  By \eqref{eq:source-16}, replacing $f$ by
$F_r$ in this pairing costs $O(n^{-\delta})+o(1)$: the feature has
norm one and
$Q_\mu|f-F_r|^2\leq Cr^{-2}n^{-2/3}+o(1)$.
The remaining pairing uses the conditional mean
$M_r^{[2]}=\mathbb E[XX^t/n\mid h_r]$.  If $X,X'$ are independent posterior
replicas, then
\[
 |M_r^{[2]}|_{\rm HS}^2
 =\mathbb E\left[\left.\left(\frac{X\cdot X'}n\right)^2
                         \right|h_r\right].
\]
The pair-tail estimate following \eqref{eq:source-9}, the annulus
estimate \eqref{eq:source-7}, and the biased-to-quenched transfer used
for \eqref{eq:source-11} give
\[
 \mathbf E|M_r^{[2]}|_{\rm HS}^2
 \leq Cr^2+C\exp(-n^\xi),
 \qquad
 Q_\mu|M_r^{[2]}|_{\rm HS}^2=o_{\mathbb P}(1).
\]
For the first inequality, split posterior overlaps at $C'r$, and on
the complementary event use the likelihood-threshold argument from
\eqref{eq:source-11}.  The second follows by $L_\circ$ comparison and
Markov's inequality.  The $n^2$ matrix components cause no loss in
\eqref{eq:source-16}, because their exponentially small remainders
remain negligible after summation.
The posterior Cauchy--Schwarz bound now makes the pairing vanish
uniformly.  As in the preceding proof, whenever the trace exceeds a
fixed $\epsilon$, the normalized resolvent optimizer has a fixed
polynomial supremum bound; in this case even a constant bound suffices.
The variational formula proves \eqref{eq:source-43}.
\end{proof}
